\begin{cabstract}
由于实测合成孔径雷达（Synthetic Aperture Radar，SAR）图像获取成本高、周期长，且在复杂任务场景下往往难以获得充足的高质量样本，利用仿真SAR图像训练目标识别模型已成为一种重要替代方案。然而，仿真SAR图像与实测SAR图像在背景噪声、目标纹理和统计分布等方面均不可避免地存在显著差异，即仿真域与实测域之间存在域偏移，导致基于仿真数据训练的模型难以直接适应实测场景。为此，本文以仿真数据域为源域，实测数据域为目标域，围绕面向仿真-实测迁移的跨域SAR目标识别方法开展研究。重点针对“仿真与少量实测数据协同训练”和“纯仿真数据训练”两类典型场景分别开展创新设计与技术验证，并进一步实现了基于国产AI芯片的模型部署。主要工作如下：

（1）针对仿真与少量实测数据协同训练场景下目标域信息不足、跨域迁移困难的问题，提出了一种基于像素级域对齐和类混淆损失的域适应SAR目标识别方法。该方法在像素级上通过图像块对比学习降低风格迁移模型对目标域样本的依赖，使仿真图像在少量实测样本条件下仍能稳定迁移到实测风格；在特征级上结合类混淆损失与一致性约束，直接优化少量实测样本条件下的目标域判别边界，减少对目标域整体分布估计的依赖。实验结果表明，所提方法有效缓解了现有域适应方法在实测样本不足条件下仿真域与实测域难以稳定对齐、模型跨域识别能力不足的问题。

（2）针对纯仿真数据训练场景下训练阶段目标域完全不可见、模型易过拟合于仿真域的问题，提出了一种基于多层级数据增强和对抗学习的域泛化SAR目标识别方法。该方法在像素级上结合高斯噪声注入、高频扰动和分贝化变换，从噪声、散射风格和动态范围等维度扩展仿真数据分布；在特征级上引入风格混淆模块对实例级特征统计量进行随机混合，模拟潜在域风格变化并抑制模型对仿真数据特征的过拟合，进一步通过域对抗学习约束模型提取域不变判别特征。实验结果表明，所提方法显著提升了基于纯仿真数据训练的SAR目标识别模型面向未见实测数据的识别性能。

（3）针对跨域SAR目标识别模型在国产边缘平台部署过程中参数规模较大、推理开销较高的问题，提出了一种面向华为昇腾NPU的轻量化部署方案。该方案结合知识蒸馏与INT8后训练量化，实现了跨域识别模型的压缩与边缘设备部署验证。实验结果表明，该方案在较好保持跨域识别性能的同时，能够有效降低模型复杂度与推理开销，为跨域SAR目标识别模型在国产边缘设备上的实际部署应用提供了支撑。

\end{cabstract}
\ckeywords{合成孔径雷达；自动目标识别；仿真-实测迁移；域适应；域泛化；知识蒸馏；模型量化}

\begin{eabstract}
Due to the high cost and lengthy acquisition cycle of measured Synthetic Aperture Radar (SAR) images, as well as the difficulty of obtaining sufficient high-quality samples under complex task scenarios, training target recognition models with simulated SAR images has emerged as an important alternative. With advances in computer-aided design and electromagnetic simulation technologies, large-scale generation of accurately labeled simulated SAR images has become feasible. However, significant discrepancies inevitably exist between simulated and measured SAR images in terms of background environment, target texture, and statistical distribution, i.e., domain shift between the simulated and measured domains, which makes it challenging for models trained on simulated data to adapt directly to measured scenarios. To this end, this thesis designates the simulated data domain as the source domain and the measured data domain as the target domain, and investigates cross-domain SAR target recognition methods oriented toward simulated-to-measured transfer. The research focuses on two representative scenarios---collaborative training with simulated and a small number of measured samples, and training with purely simulated data---conducting innovative method design and technical validation for each, and further realizes model deployment on a domestic AI chip. The main contributions are as follows:

(1) To address the issues of insufficient target-domain information and difficult inter-domain alignment under the scenario of collaborative training with simulated and a small number of measured samples, a domain adaptation method for SAR target recognition based on pixel-level domain alignment and class confusion loss is proposed. The method alleviates the cross-domain discrepancy between simulated and measured domains through pixel-level domain alignment, and enhances target-domain discrimination under limited measured samples by incorporating class confusion loss. Experimental results demonstrate that the proposed method achieves effective simulated-to-measured domain transfer with only a small number of measured samples, effectively mitigating the difficulty of stable inter-domain alignment and insufficient cross-domain recognition capability faced by existing domain adaptation methods under scarce measured data conditions.

(2) To address the issues that the target domain is entirely unavailable during training and the model is prone to overfitting the simulated domain under the purely simulated data training scenario, a domain generalization method for SAR target recognition based on multi-level data augmentation and adversarial learning is proposed. The method expands the simulated data distribution at the pixel level through Gaussian noise injection, high-frequency perturbation, and decibel transformation, introduces the MixStyle module at the feature level to randomly mix instance-wise statistics and simulate latent domain style variations, and leverages domain adversarial learning to constrain the extraction of domain-invariant discriminative features. Experimental results demonstrate that the proposed method effectively enhances the model's generalization capability and cross-domain stability on unseen measured data, and yields superior cross-domain recognition performance to the baseline and representative comparison methods under the purely simulated training setting.

(3) To address the issues of large parameter scales and high inference overhead encountered when deploying cross-domain SAR target recognition models on domestic edge platforms, a lightweight deployment scheme targeting Huawei Ascend chips is proposed. The scheme integrates knowledge distillation with INT8 post-training quantization to achieve compression and on-board deployment validation of cross-domain recognition models. Experimental results demonstrate that the proposed scheme effectively reduces model complexity and inference overhead while largely preserving cross-domain recognition performance, thereby providing support for the practical deployment of cross-domain SAR target recognition models on domestic edge devices.

\end{eabstract}
\ekeywords{Synthetic Aperture Radar, Automatic Target Recognition, Cross-Domain Transfer Learning, Domain Adaptation, Domain Generalization, Knowledge Distillation, Model Quantization}
